\sourcepage{164}
\clearpage
\chapterappendix{6A}{用 \texorpdfstring{$t$}{t} 检验比较连续变量均值时，每个等大组的所需样本量}
\sourcepage{165}
\subsection*{表 6A\quad 比较两个均值的每组样本量}
最左列为标准化效应量 $E/S$。
{\small\setlength{\tabcolsep}{4pt}
\begin{longtable}{lrrrrrrrrr}
\toprule
单侧 $\alpha$ & \multicolumn{3}{c}{0.005} & \multicolumn{3}{c}{0.025} & \multicolumn{3}{c}{0.05}\\
双侧 $\alpha$ & \multicolumn{3}{c}{0.01} & \multicolumn{3}{c}{0.05} & \multicolumn{3}{c}{0.10}\\
$\beta$ & 0.05 & 0.10 & 0.20 & 0.05 & 0.10 & 0.20 & 0.05 & 0.10 & 0.20\\\midrule
\endfirsthead
\toprule
单侧 $\alpha$ & \multicolumn{3}{c}{0.005} & \multicolumn{3}{c}{0.025} & \multicolumn{3}{c}{0.05}\\
双侧 $\alpha$ & \multicolumn{3}{c}{0.01} & \multicolumn{3}{c}{0.05} & \multicolumn{3}{c}{0.10}\\
$\beta$ & 0.05 & 0.10 & 0.20 & 0.05 & 0.10 & 0.20 & 0.05 & 0.10 & 0.20\\\midrule
\endhead
0.10 & 3565 & 2978 & 2338 & 2600 & 2103 & 1571 & 2166 & 1714 & 1238\\
0.15 & 1586 & 1325 & 1040 & 1157 & 935 & 699 & 963 & 762 & 551\\
0.20 & 893 & 746 & 586 & 651 & 527 & 394 & 542 & 429 & 310\\
0.25 & 572 & 478 & 376 & 417 & 338 & 253 & 347 & 275 & 199\\
0.30 & 398 & 333 & 262 & 290 & 235 & 176 & 242 & 191 & 139\\
0.40 & 225 & 188 & 148 & 164 & 133 & 100 & 136 & 108 & 78\\
0.50 & 145 & 121 & 96 & 105 & 86 & 64 & 88 & 70 & 51\\
0.60 & 101 & 85 & 67 & 74 & 60 & 45 & 61 & 49 & 36\\
0.70 & 75 & 63 & 50 & 55 & 44 & 34 & 45 & 36 & 26\\
0.80 & 58 & 49 & 39 & 42 & 34 & 26 & 35 & 28 & 21\\
0.90 & 46 & 39 & 32 & 34 & 27 & 21 & 28 & 22 & 16\\
1.00 & 38 & 32 & 26 & 27 & 23 & 17 & 23 & 18 & 14\\\bottomrule
\end{longtable}}

\sourcepage{166}
\noindent\small $E/S$ 为预期效应量 $E$（两组均值差）除以结局标准差 $S$。按该行与指定 alpha、beta 列查得每组人数。单样本 $t$ 检验的总人数，为表列数字的一半。\par\normalsize
\subsection{计算变异程度}
变异常报告为标准差，或均值标准误（standard error of the mean，SEM）。样本量计算最有用的是标准差，转换很容易：标准差等于标准误乘以 $\sqrt N$，其中 $N$ 为用于计算平均值的参与者人数。例如，25 人低纤维饮食后减重为 $10\pm2$ kg（均值 $\pm$ SEM），则标准差为 $2\text{ kg}\times\sqrt{25}=10\text{ kg}$。
\subsection{其他取值}
表中没有的 $E/S$，或两组大小不等时，请使用 \url{https://www.sample-size.net} 计算器。

\clearpage
\chapterappendix{6B}{用卡方统计量或 \texorpdfstring{$Z$}{Z} 检验比较二分类变量比例时，每个等大组的所需样本量}
\sourcepage{167}
\subsection*{表 6B.1\quad 比较两个比例的每组样本量}
每格三数分别对应：上方，单侧 $\alpha=0.05$ 或双侧 $\alpha=0.10$，$\beta=0.20$；中间，单侧 $\alpha=0.025$ 或双侧 $\alpha=0.05$，$\beta=0.20$；下方，单侧 $\alpha=0.025$ 或双侧 $\alpha=0.05$，$\beta=0.10$。各数值列为 $|P_1-P_0|$，最左列为两比例中较小者。
{\small\setlength{\tabcolsep}{3.5pt}
\begin{longtable}{lrrrrrrrrrr}
\toprule $\min(P_0,P_1)$ & 0.05 & 0.10 & 0.15 & 0.20 & 0.25 & 0.30 & 0.35 & 0.40 & 0.45 & 0.50\\\midrule
\endfirsthead
\toprule $\min(P_0,P_1)$ & 0.05 & 0.10 & 0.15 & 0.20 & 0.25 & 0.30 & 0.35 & 0.40 & 0.45 & 0.50\\\midrule
\endhead
0.05 & 381 & 129 & 72 & 47 & 35 & 27 & 22 & 18 & 15 & 13\\*
 & 473 & 159 & 88 & 59 & 43 & 33 & 26 & 22 & 18 & 16\\*
 & 620 & 207 & 113 & 75 & 54 & 41 & 33 & 27 & 23 & 19\\\midrule
0.10 & 578 & 175 & 91 & 58 & 41 & 31 & 24 & 20 & 16 & 14\\*
 & 724 & 219 & 112 & 72 & 51 & 37 & 29 & 24 & 20 & 17\\*
 & 958 & 286 & 146 & 92 & 65 & 48 & 37 & 30 & 25 & 21\\\midrule
0.15 & 751 & 217 & 108 & 67 & 46 & 34 & 26 & 21 & 17 & 15\\*
 & 944 & 270 & 133 & 82 & 57 & 41 & 32 & 26 & 21 & 18\\*
 & 1252 & 354 & 174 & 106 & 73 & 53 & 42 & 33 & 26 & 22\\\midrule
0.20 & 900 & 251 & 121 & 74 & 50 & 36 & 28 & 22 & 18 & 15\\*
 & 1133 & 313 & 151 & 91 & 62 & 44 & 34 & 27 & 22 & 18\\*
 & 1504 & 412 & 197 & 118 & 80 & 57 & 44 & 34 & 27 & 23\\\midrule
0.25 & 1024 & 278 & 132 & 79 & 53 & 38 & 29 & 23 & 18 & 15\\*
 & 1289 & 348 & 165 & 98 & 66 & 47 & 35 & 28 & 22 & 18\\*
 & 1714 & 459 & 216 & 127 & 85 & 60 & 46 & 35 & 28 & 23\\\midrule
0.30 & 1123 & 300 & 141 & 83 & 55 & 39 & 29 & 23 & 18 & 15\\*
 & 1415 & 376 & 175 & 103 & 68 & 48 & 36 & 28 & 22 & 18\\*
 & 1883 & 496 & 230 & 134 & 88 & 62 & 47 & 36 & 28 & 23\\\midrule
0.35 & 1197 & 315 & 146 & 85 & 56 & 39 & 29 & 23 & 18 & 15\\*
 & 1509 & 395 & 182 & 106 & 69 & 48 & 36 & 28 & 22 & 18\\*
 & 2009 & 522 & 239 & 138 & 90 & 62 & 47 & 35 & 27 & 22\\\midrule
0.40 & 1246 & 325 & 149 & 86 & 56 & 39 & 29 & 22 & 17 & 14\\*
 & 1572 & 407 & 186 & 107 & 69 & 48 & 35 & 27 & 21 & 17\\*
 & 2093 & 538 & 244 & 139 & 90 & 62 & 46 & 34 & 26 & 21\\\midrule
0.45 & 1271 & 328 & 149 & 85 & 55 & 38 & 28 & 21 & 16 & 13\\*
 & 1603 & 411 & 186 & 106 & 68 & 47 & 34 & 26 & 20 & 16\\*
 & 2135 & 543 & 244 & 138 & 88 & 60 & 44 & 33 & 25 & 19\\\midrule
0.50 & 1271 & 325 & 146 & 83 & 53 & 36 & 26 & 20 & 15 & —\\*
 & 1603 & 407 & 182 & 103 & 66 & 44 & 32 & 24 & 18 & —\\*
 & 2135 & 538 & 239 & 134 & 85 & 57 & 42 & 30 & 23 & —\\\midrule
0.55 & 1246 & 315 & 141 & 79 & 50 & 34 & 24 & 18 & — & —\\*
 & 1572 & 395 & 175 & 98 & 62 & 41 & 29 & 22 & — & —\\*
 & 2093 & 522 & 230 & 127 & 80 & 53 & 37 & 27 & — & —\\\midrule
0.60 & 1197 & 300 & 132 & 74 & 46 & 31 & 22 & — & — & —\\*
 & 1509 & 376 & 165 & 91 & 57 & 37 & 26 & — & — & —\\*
 & 2009 & 496 & 216 & 118 & 73 & 48 & 33 & — & — & —\\\midrule
0.65 & 1123 & 278 & 121 & 67 & 41 & 27 & — & — & — & —\\*
 & 1415 & 348 & 151 & 82 & 51 & 33 & — & — & — & —\\*
 & 1883 & 459 & 197 & 106 & 65 & 41 & — & — & — & —\\\midrule
0.70 & 1024 & 251 & 108 & 58 & 35 & — & — & — & — & —\\*
 & 1289 & 313 & 133 & 72 & 43 & — & — & — & — & —\\*
 & 1714 & 412 & 174 & 92 & 54 & — & — & — & — & —\\\midrule
0.75 & 900 & 217 & 91 & 47 & — & — & — & — & — & —\\*
 & 1133 & 270 & 112 & 59 & — & — & — & — & — & —\\*
 & 1504 & 354 & 146 & 75 & — & — & — & — & — & —\\\midrule
0.80 & 751 & 175 & 72 & — & — & — & — & — & — & —\\*
 & 944 & 219 & 88 & — & — & — & — & — & — & —\\*
 & 1252 & 286 & 113 & — & — & — & — & — & — & —\\\midrule
0.85 & 578 & 129 & — & — & — & — & — & — & — & —\\*
 & 724 & 159 & — & — & — & — & — & — & — & —\\*
 & 958 & 207 & — & — & — & — & — & — & — & —\\\midrule
0.90 & 381 & — & — & — & — & — & — & — & — & —\\*
 & 473 & — & — & — & — & — & — & — & — & —\\*
 & 620 & — & — & — & — & — & — & — & — & —\\\midrule
\end{longtable}}

\sourcepage{168}
单侧估计使用 $Z$ 统计量。$P_0$ 表示一组中预期发生结局的比例，$P_1$ 是另一组的比例。在病例对照研究，$P_1$ 是病例中具有预测特征者比例，$P_0$ 是对照中相应比例。按较小的 $P_0$ 或 $P_1$ 行和预期绝对差列读取；三个数表示指定 alpha、beta 下各组所需样本量。$P_0$ 或 $P_1$ 为 0.01～0.10 的更详细数字见表 6B.2。
\subsection{其他取值}
\sourcepage{169}
两表未列的比例，或两组大小不等时，请用网站计算器。
\subsection*{表 6B.2\quad 比较两个比例的每组样本量：较小比例为 0.01～0.10}
上、中、下三数的 alpha、beta 组合与表 6B.1 相同；各列为两比例绝对差，最左列为较小比例。
{\small\setlength{\tabcolsep}{3.5pt}
\begin{longtable}{lrrrrrrrrrr}
\toprule $\min(P_0,P_1)$ & 0.01 & 0.02 & 0.03 & 0.04 & 0.05 & 0.06 & 0.07 & 0.08 & 0.09 & 0.10\\\midrule
\endfirsthead
\toprule $\min(P_0,P_1)$ & 0.01 & 0.02 & 0.03 & 0.04 & 0.05 & 0.06 & 0.07 & 0.08 & 0.09 & 0.10\\\midrule
\endhead
0.01 & 2019 & 700 & 396 & 271 & 204 & 162 & 134 & 114 & 98 & 87\\*
 & 2512 & 864 & 487 & 332 & 249 & 197 & 163 & 138 & 120 & 106\\*
 & 3300 & 1125 & 631 & 428 & 320 & 254 & 209 & 178 & 154 & 135\\\midrule
0.02 & 3205 & 994 & 526 & 343 & 249 & 193 & 157 & 131 & 113 & 97\\*
 & 4018 & 1237 & 651 & 423 & 306 & 238 & 192 & 161 & 137 & 120\\*
 & 5320 & 1625 & 852 & 550 & 397 & 307 & 248 & 207 & 177 & 154\\\midrule
0.03 & 4367 & 1283 & 653 & 414 & 294 & 224 & 179 & 148 & 126 & 109\\*
 & 5493 & 1602 & 813 & 512 & 363 & 276 & 220 & 182 & 154 & 133\\*
 & 7296 & 2114 & 1067 & 671 & 474 & 359 & 286 & 236 & 199 & 172\\\midrule
0.04 & 5505 & 1564 & 777 & 482 & 337 & 254 & 201 & 165 & 139 & 119\\*
 & 6935 & 1959 & 969 & 600 & 419 & 314 & 248 & 203 & 170 & 146\\*
 & 9230 & 2593 & 1277 & 788 & 548 & 410 & 323 & 264 & 221 & 189\\\midrule
0.05 & 6616 & 1838 & 898 & 549 & 380 & 283 & 222 & 181 & 151 & 129\\*
 & 8347 & 2308 & 1123 & 686 & 473 & 351 & 275 & 223 & 186 & 159\\*
 & 11123 & 3061 & 1482 & 902 & 620 & 460 & 360 & 291 & 242 & 206\\\midrule
0.06 & 7703 & 2107 & 1016 & 615 & 422 & 312 & 243 & 197 & 163 & 139\\*
 & 9726 & 2650 & 1272 & 769 & 526 & 388 & 301 & 243 & 202 & 171\\*
 & 12973 & 3518 & 1684 & 1014 & 691 & 508 & 395 & 318 & 263 & 223\\\midrule
0.07 & 8765 & 2369 & 1131 & 680 & 463 & 340 & 263 & 212 & 175 & 148\\*
 & 11076 & 2983 & 1419 & 850 & 577 & 423 & 327 & 263 & 217 & 183\\*
 & 14780 & 3965 & 1880 & 1123 & 760 & 555 & 429 & 343 & 283 & 239\\\midrule
0.08 & 9803 & 2627 & 1244 & 743 & 502 & 367 & 282 & 227 & 187 & 158\\*
 & 12393 & 3308 & 1562 & 930 & 627 & 457 & 352 & 282 & 232 & 195\\*
 & 16546 & 4401 & 2072 & 1229 & 827 & 602 & 463 & 369 & 303 & 255\\\midrule
0.09 & 10816 & 2877 & 1354 & 804 & 541 & 393 & 302 & 241 & 198 & 167\\*
 & 13679 & 3626 & 1702 & 1007 & 676 & 491 & 377 & 300 & 246 & 207\\*
 & 18270 & 4827 & 2259 & 1333 & 893 & 647 & 495 & 393 & 322 & 270\\\midrule
0.10 & 11804 & 3121 & 1461 & 863 & 578 & 419 & 320 & 255 & 209 & 175\\*
 & 14933 & 3936 & 1838 & 1083 & 724 & 523 & 401 & 318 & 260 & 218\\*
 & 19952 & 5242 & 2441 & 1434 & 957 & 690 & 527 & 417 & 341 & 285\\\midrule
\end{longtable}}

\sourcepage{170}
单侧估计使用 $Z$ 统计量。
\subsection{其他取值}
两表未列的 $P_0$、$P_1$，或两组大小不等时，请使用 \url{https://www.sample-size.net}。

\clearpage
\chapterappendix{6C}{使用相关系数时所需的总样本量}
\sourcepage{171}
\subsection*{表 6C\quad 判断相关系数是否异于零的样本量}
最左列为预期相关系数 $r$。按该行和指定 alpha、beta 列查总人数。
{\small\setlength{\tabcolsep}{4pt}
\begin{longtable}{lrrrrrrrrr}
\toprule
单侧 $\alpha$ & \multicolumn{3}{c}{0.005} & \multicolumn{3}{c}{0.025} & \multicolumn{3}{c}{0.05}\\
双侧 $\alpha$ & \multicolumn{3}{c}{0.01} & \multicolumn{3}{c}{0.05} & \multicolumn{3}{c}{0.10}\\
$\beta$ & 0.05 & 0.10 & 0.20 & 0.05 & 0.10 & 0.20 & 0.05 & 0.10 & 0.20\\\midrule
\endfirsthead
\toprule
单侧 $\alpha$ & \multicolumn{3}{c}{0.005} & \multicolumn{3}{c}{0.025} & \multicolumn{3}{c}{0.05}\\
双侧 $\alpha$ & \multicolumn{3}{c}{0.01} & \multicolumn{3}{c}{0.05} & \multicolumn{3}{c}{0.10}\\
$\beta$ & 0.05 & 0.10 & 0.20 & 0.05 & 0.10 & 0.20 & 0.05 & 0.10 & 0.20\\\midrule
\endhead
0.05 & 7118 & 5947 & 4663 & 5193 & 4200 & 3134 & 4325 & 3424 & 2469\\
0.10 & 1773 & 1481 & 1162 & 1294 & 1047 & 782 & 1078 & 854 & 616\\
0.15 & 783 & 655 & 514 & 572 & 463 & 346 & 477 & 378 & 273\\
0.20 & 436 & 365 & 287 & 319 & 259 & 194 & 266 & 211 & 153\\
0.25 & 276 & 231 & 182 & 202 & 164 & 123 & 169 & 134 & 98\\
0.30 & 189 & 158 & 125 & 139 & 113 & 85 & 116 & 92 & 67\\
0.35 & 136 & 114 & 90 & 100 & 82 & 62 & 84 & 67 & 49\\
0.40 & 102 & 86 & 68 & 75 & 62 & 47 & 63 & 51 & 37\\
0.45 & 79 & 66 & 53 & 58 & 48 & 36 & 49 & 39 & 29\\
0.50 & 62 & 52 & 42 & 46 & 38 & 29 & 39 & 31 & 23\\
0.60 & 40 & 34 & 27 & 30 & 25 & 19 & 26 & 21 & 16\\
0.70 & 27 & 23 & 19 & 20 & 17 & 13 & 17 & 14 & 11\\
0.80 & 18 & 15 & 13 & 14 & 12 & 9 & 12 & 10 & 8\\\bottomrule
\end{longtable}}

\subsection{其他取值}
表中未列的 $r$，请使用网站计算器。

\sourcepage{172}
\clearpage
\chapterappendix{6D}{连续变量描述性研究的样本量}
\sourcepage{173}
\subsection*{表 6D\quad 常见 \texorpdfstring{$W/S$}{W/S} 下的样本量}
数字列为置信水平。
{\small\setlength{\tabcolsep}{4pt}
\begin{longtable}{lrrr}
\toprule $W/S$ & 90\% & 95\% & 99\%\\\midrule
\endfirsthead
\toprule $W/S$ & 90\% & 95\% & 99\%\\\midrule
\endhead
0.10 & 1083 & 1537 & 2665\\
0.15 & 482 & 683 & 1180\\
0.20 & 271 & 385 & 664\\
0.25 & 174 & 246 & 425\\
0.30 & 121 & 171 & 295\\
0.35 & 89 & 126 & 217\\
0.40 & 68 & 97 & 166\\
0.50 & 44 & 62 & 107\\
0.60 & 31 & 43 & 74\\
0.70 & 23 & 32 & 55\\
0.80 & 17 & 25 & 42\\
0.90 & 14 & 19 & 33\\
1.00 & 11 & 16 & 27\\\bottomrule
\end{longtable}}

$W/S$ 是标准化置信区间宽度，即所需总宽 $W$ 除以变量标准差 $S$。按标准化宽度行与指定置信水平列查总人数。
\subsection{其他取值}
表中未列的 $W/S$，请使用网站计算器。

\clearpage
\chapterappendix{6E}{二分类变量描述性研究的样本量}
\sourcepage{174}
\subsection*{表 6E\quad 比例的样本量}
最左列为预期比例 $P$，各数值列为置信区间总宽 $W$。每格上、中、下三数依次为 90\%、95\%、99\% 置信水平所需样本量。
{\small\setlength{\tabcolsep}{3.5pt}
\begin{longtable}{lrrrrrrr}
\toprule $P$ & 0.10 & 0.15 & 0.20 & 0.25 & 0.30 & 0.35 & 0.40\\\midrule
\endfirsthead
\toprule $P$ & 0.10 & 0.15 & 0.20 & 0.25 & 0.30 & 0.35 & 0.40\\\midrule
\endhead
0.10 & 98 & 44 & — & — & — & — & —\\*
 & 138 & 61 & — & — & — & — & —\\*
 & 239 & 106 & — & — & — & — & —\\\midrule
0.15 & 139 & 62 & 35 & 22 & — & — & —\\*
 & 196 & 87 & 49 & 31 & — & — & —\\*
 & 339 & 151 & 85 & 54 & — & — & —\\\midrule
0.20 & 174 & 77 & 44 & 28 & 19 & 14 & —\\*
 & 246 & 109 & 61 & 39 & 27 & 20 & —\\*
 & 426 & 189 & 107 & 68 & 47 & 35 & —\\\midrule
0.25 & 204 & 91 & 51 & 33 & 23 & 17 & 13\\*
 & 288 & 128 & 72 & 46 & 32 & 24 & 18\\*
 & 499 & 222 & 125 & 80 & 55 & 41 & 31\\\midrule
0.30 & 229 & 102 & 57 & 37 & 25 & 19 & 14\\*
 & 323 & 143 & 81 & 52 & 36 & 26 & 20\\*
 & 559 & 249 & 140 & 89 & 62 & 46 & 35\\\midrule
0.40 & 261 & 116 & 65 & 42 & 29 & 21 & 16\\*
 & 369 & 164 & 92 & 59 & 41 & 30 & 23\\*
 & 639 & 284 & 160 & 102 & 71 & 52 & 40\\\midrule
0.50 & 272 & 121 & 68 & 44 & 30 & 22 & 17\\*
 & 384 & 171 & 96 & 61 & 43 & 31 & 24\\*
 & 666 & 296 & 166 & 107 & 74 & 54 & 42\\\midrule
\end{longtable}}

\sourcepage{175}
按具有目标特征的预期比例行和所需区间总宽列查人数；三数对应三个置信水平。
\subsection{其他取值}
表中未列的 $P$，请使用网站计算器。

\clearpage
\chapterappendix{6F}{\texorpdfstring{$t$}{t} 检验的使用与误用}
\sourcepage{176}
本章重点之一的双样本 $t$ 检验，比较两组变量均值。组别可按二分类预测变量定义，例如随机试验的有效药与安慰剂、队列的危险因素有无；也可按二分类结局定义，例如病例对照研究。比较一次测量时为非配对；比较每人两个时间点，例如干预前后测量的变化时，本书称配对。第三种单样本配对 $t$ 检验，则比较单组个体两次测量的平均变化与零或其他指定变化值。

表 6F 展示了原本应作组间比较，却误用单样本配对检验的情况：随机盲法试验研究新安眠药对 0～10 分生活质量的影响。一些研究者曾分别对治疗及安慰剂组各做一次单样本检验，甚至发表，这并不恰当。

标有匕首号 $\dagger$ 的 $P$ 值来自单样本配对检验。治疗组 $P=0.05$，显示研究期间显著变化；对照组 $P=0.16$，无显著变化。但这不能推断组间差别，不能据此说治疗有显著效果。

星号 $*$ 是适当的双样本检验。前两个 $P=0.87$ 和 0.64 为非配对，显示初始及末次生活质量组间均无显著差异。最后 $P=0.17$ 是本书所称双样本配对检验，比末次比较 0.64 更接近 0.05，因为变化均值的标准差更小。
\sourcepage{177}
但是治疗组平均改善 1.3 分，与安慰剂 0.9 分无显著差异。因此，正确结论是本研究未发现治疗有效。
\begin{bookbox}{表 6F\quad 分析配对资料的正确与错误方式}
\begin{tabularx}{\linewidth}{YYYY}
\toprule 测量时间 & 治疗（$N=100$） & 对照（$N=100$） & $P$ 值\\
\multicolumn{4}{c}{生活质量：均值 $\pm$ 标准差}\\\midrule
研究开始 & $7.0\pm4.5$ & $7.1\pm4.4$ & $0.87^*$\\
研究结束 & $8.3\pm4.7$ & $8.0\pm4.6$ & $0.64^*$\\
变化（结束减开始） & $1.3\pm2.1$ & $0.9\pm2.0$ & $0.17^*$\\
$P$ 值 & $0.05^\dagger$ & $0.16^\dagger$ & \\\bottomrule
\end{tabularx}
\footnotesize $*$ 比较治疗组与对照组。$\dagger$ 比较变化值与零。
\end{bookbox}

\clearpage
\chapterappendix{6G}{第 6 章习题：估计样本量与效能}
\sourcepage{178}
\begin{enumerate}
\item 回顾第 5 章习题 1。研究需多少胃癌病例？若效能要求 0.90 呢？若显著性水平 0.01 呢？附加题：如果只有 60 例，她可以怎样处理？
\item 肌力随年龄增长下降。初步资料提示，部分可能由脱氢表雄酮（dehydroepiandrosterone，DHEA）逐渐缺乏造成。研究者计划随机给予老年人 DHEA 或相同外观安慰剂六个月，再测肌力。既往平均握力 20 kg、标准差 8 kg。双侧 alpha 0.05、beta 0.10 时，为显示治疗组与安慰剂至少 10\% 差异，需多少人？若 beta 0.20 呢？
\item 习题 2 计算人数超过可招募数。同事指出，老年人握力差异很大，可能掩盖治疗效果；建议研究开始及治疗后各测一次，以变化作结局。小型先导研究显示六个月变化标准差仅 2 kg。双侧 alpha 0.05、beta 0.10，仍检出效应量 2 kg，每组需多少人？
\item 研究者怀疑阅读障碍三年级学生中左利手更多。既往约 10\% 为左利手，阅读障碍不常见。计划选学区全部阅读障碍者为病例，随机选同样多无障碍者为对照。若检出左利手相对右利手的阅读障碍 OR 为 2.0，双侧 alpha 0.05、beta 0.20，需要多少人？
\item 研究者希望估计所在机构医学生平均智商，99\% 置信区间为 $\pm3$ 分。小型先导研究显示约 110～150 分，需要的大致人数是多少？
\end{enumerate}
